% Draft (21 Oct 2026) of the new results for the coupling of Section 4.6, to replace the paragraph "Two species"
% and Figures fig:ch5-map, fig:ch5-spreading in ch4_ansys.tex. Source of the set-up: MEETING_2026-10-20.md.
% ANSYS numbers: 8^3, the fixed build (A-22/A-23, runs v2f_*, fsens_*); [16^3] to fill from A-23/A-24.
% Python: coupled_living_consumption.py --v2 --phi-min 0.1 --front 0.1 --fortran --ansys-bc --start-position.
% Not in the build yet.
% 21 Oct, A-25: every ANSYS number below (v2f_*, fsens_*) comes from runs in which phi advanced (outer/inner) = 10
% times too fast (inner USSFin step 0.0025). They are to be replaced by the u1_ runs (one inner step, as in Chapter 4).
% The "outer increment halved" row and its paragraph are this error, not a time-step convergence issue; the
% replay, hand-over and constraint checks stay valid. The Python numbers are not affected.

\paragraph{The coupled point model against the point model alone.}
With uniform fields, $c=1$ everywhere and the antibiotic $a=1$ from the start, every point model in the
coupled run must follow the point model alone. Fritsch et al.'s synthetic case~I
\cite{Fritsch2025BayesianMicrofilms} run this way agrees with the point model alone to $2\times10^{-14}$
(Python, with the Fortran point model of the ANSYS runs and with an independent implementation in JAX,
which agree with each other to $2\times10^{-16}$). In ANSYS every step of the traced Gauss points is
recomputed from its start state and its inputs $(c_n,c_{n+1},a_n,a_{n+1})$; they agree to
$1.8\times10^{-12}$ in every step, once Newton's method is iterated to its tolerance: the first step from
the start state needs up to 38 iterations, a fixed budget of 6 or 12 leaves it unconverged. The constraint
$\sum_i\phi_i+\phi_0=1$ of the point model holds to $2\times10^{-16}$ at every traced step.
The hand-over was checked separately, because the replay above sees only the point model: the living volume
fraction that the field equations use equals the one the point model returned in the same increment, at every
increment (difference $0$), and the start state of each step is the end state of the step before.

\paragraph{Two examples.}
Following the partner's suggestion, two examples use the same model with different start values
(synthetic data, case~I of Fritsch et al., two species, $b=(1,2)$): species A in the seed of this chapter
and species B in four small blocks at the corners, and both species mixed in the seed in equal amounts
($\phi_1=\phi_2=0.275$, the total of Fritsch et al.). The nutrient is held on the layer $y=-1$~mm, the
antibiotic on all outer faces from $T^*=0.5$, with $\alpha^*_0=50$ and a weaker uptake than the nutrient
($g_aL^2/d_a=0.25$ against $4$; an assumption). The front term uses a tenth of the converted Table~2 value
of $r$; with the full value the elements collapse within the first increments. A point starts its point model
when $\phi_{3D}$ first exceeds $\phi_{\min}=0.1$. The papers do not say from which composition it starts
(Section~\ref{sec:ch5-outlook-python}); here it starts as species B within $0.5625$~mm of the vertical edges
and as species A elsewhere in the first example, and from the mixed start in the second. This is an
assumption: it fixes where the two regions meet. Starting instead from the mean state of the colonised
neighbours changes the living volume fraction of the seed by $0.025$ (Python).
Table~\ref{tab:ch5-living} gives the seed and the corner blocks at $T^*=1$.

\begin{table}[htbp]
  \centering
  \caption[The coupled examples]{%
    The two examples in ANSYS ($8^3$ [$16^3$]) at $T^*=1$, means over the seed and over the corner blocks:
    biofilm $\phi_{3D}$, living volume fraction $\sum_i\bar\phi_i$, share of species A, nutrient $c$,
    antibiotic $a$, growth $\alpha-1$ and von Mises stress.}
  \label{tab:ch5-living}
  \begin{tabular}{llccccccc}
    \toprule
    & & $\phi_{3D}$ & $\sum_i\bar\phi_i$ & A & $c$ & $a$ & $\alpha-1$ & $\sigma_{vM}$ [Pa]\\
    \midrule
    separate, antibiotic & seed & 0.579 & 0.567 & 0.969 & 0.475 & 0.990 & $6.35\cdot10^{-4}$ & $1.57\cdot10^{-4}$\\
    & corners & 0.614 & 0.448 & 0.031 & 0.688 & 1.000 & $5.40\cdot10^{-4}$ & $1.61\cdot10^{-4}$\\
    separate & seed & [..] & & & & & & \\
    mixed, antibiotic & seed & 0.542 & 0.075 & 0.641 & 0.849 & 0.998 & $7.00\cdot10^{-4}$ & $1.24\cdot10^{-4}$\\
    mixed & seed & 0.513 & 0.949 & 0.688 & 0.376 & 0 & $6.95\cdot10^{-4}$ & $1.11\cdot10^{-4}$\\
    \bottomrule
  \end{tabular}
\end{table}

\paragraph{What the antibiotic does.}
Species B, the more sensitive one ($b_2=2$), dies wherever the antibiotic arrives. In the mixed seed the
living volume fraction falls from $0.95$ to $0.08$, and the nutrient, no longer consumed, rises from $0.38$
to $0.85$. Whether a species survives depends on the balance in its equation for $\psi_i$ between the
nutrient term $c^*\phi_i\sum_jA_{ij}\phi_j\psi_j$ and the antibiotic term $\alpha^*b_i\psi_i$: with
$c^*=100\,c$ and $\alpha^*b_1=50$, species A grows where $c>0.5$ and declines where $c<0.5$. The separate seed
lies near this threshold ($c=0.475$), and its living volume fraction at $T^*=1$ follows $\alpha^*_0$ in
steps: $0.94$, $0.73$, $0.57$ and $0.42$ for $\alpha^*_0=10$, $40$, $50$ and $60$. The nutrient field thus
decides where the antibiotic acts, which a point model with a constant $c^*$ cannot show.

The mixed example also depends on the start: with the start of Fritsch et al.\ ($\phi_1=0.25$,
$\phi_2=0.30$) the share of A in the seed ends at $0.52$ and the living volume fraction at $0.048$, with
equal amounts at $0.64$ and $0.075$ (ANSYS; Python $0.57$/$0.050$ and $0.77$/$0.21$). As in case~6 of
Section~\ref{sec:ch5-outlook-python}, the start share decides which species prevails.

The growth $\alpha-1$ and the stress hardly change with the antibiotic ($6.95$ against
$7.00\cdot10^{-4}$ in the mixed seed). This follows from the coupling: Eq.~36 is unchanged, so the
biofilm expands with $\phi_{3D}$ whether it is alive or dead, and the point model reaches the stress only
through the consumption, the nutrient field and the front term of Eq.~34.

\paragraph{Python against ANSYS.}
The Python version with the boundary conditions of the ANSYS deck (the nutrient and the antibiotic held on
whole element layers) and the same start rule gives in the separate seed at $T^*=1$ a living fraction of
$0.830$ and $c=0.543$ (ANSYS $0.567$ and $0.475$), in the mixed seed with equal amounts $0.207$ and $0.851$
(ANSYS $0.075$ and $0.849$). The trends agree; the values differ most where the seed lies near the threshold
$c=0.5$. In both the seed's $\phi_{3D}$ falls below 1 as the front term moves biomass towards the nutrient.
The differences come from the discretisation of the fields ($32^3$ finite volumes against $8^3$ elements
of the neighboured element method, with the time step discussed below).

\paragraph{Sensitivity.}
Table~\ref{tab:ch5-living-sens} varies one choice at a time in the separate example with antibiotic.

\begin{table}[htbp]
  \centering
  \caption[Sensitivity of the coupled example]{%
    Separate example with antibiotic, ANSYS $8^3$, seed at $T^*=1$: one choice changed at a time.}
  \label{tab:ch5-living-sens}
  \begin{tabular}{lcccccc}
    \toprule
    & $\phi_{3D}$ & $\sum_i\bar\phi_i$ & $c$ & $a$ & $\alpha-1$ & $\sigma_{vM}$ [Pa]\\
    \midrule
    reference & 0.579 & 0.567 & 0.475 & 0.990 & $6.35\cdot10^{-4}$ & $1.57\cdot10^{-4}$\\
    $\phi_{\min}=0.01$ & 0.561 & 0.572 & 0.476 & 0.990 & $6.08\cdot10^{-4}$ & $1.51\cdot10^{-4}$\\
    outer increment halved & 0.488 & 0.629 & 0.507 & 0.991 & $7.01\cdot10^{-4}$ & $2.22\cdot10^{-4}$\\
    antibiotic uptake $\times16$ & 0.576 & 0.734 & 0.453 & 0.834 & $6.35\cdot10^{-4}$ & $1.56\cdot10^{-4}$\\
    $\alpha^*_0=10$ & 0.569 & 0.937 & 0.350 & 0.984 & $6.33\cdot10^{-4}$ & $1.52\cdot10^{-4}$\\
    $\alpha^*_0=40$ & 0.576 & 0.733 & 0.422 & 0.988 & $6.34\cdot10^{-4}$ & $1.55\cdot10^{-4}$\\
    $\alpha^*_0=60$ & 0.582 & 0.418 & 0.540 & 0.992 & $6.36\cdot10^{-4}$ & $1.59\cdot10^{-4}$\\
    \bottomrule
  \end{tabular}
\end{table}

The choices of the point model ($\phi_{\min}$, $\alpha^*_0$, the antibiotic uptake) move the living volume
fraction but hardly the growth and the stress. The outer increment moves everything: halving it lowers
$\phi_{3D}$ in the seed by $16\,\%$ and raises the stress by $41\,\%$. The coupled runs are therefore not
converged in time at $\Delta T^*=0.025$. The Python version, which integrates the fields with its own finer
steps inside each outer step, changes by less than $0.5\,\%$ when the outer step is halved (living volume
fraction $0.830$ and $0.834$, stress $1.082$ and $1.083\cdot10^{-3}$~Pa on $32^3$; $0.829$ and $0.832$ on
$16^3$), so the dependence lies in the time integration of the fields in the element [ANSYS: ..].

\paragraph{A limitation of the discretisation.}
The examples are mirror-symmetric in $x$ and $z$, the results are not: on $8^3$ $\phi_{3D}$ differs by up to
$0.0085$ between mirrored elements with only the diffusion of $\phi_{3D}$, and by up to $0.09$ once the front
term amplifies it; in the runs of Section~[..] on $16^3$ by up to $0.9\,\%$ in $x$. The cause is the neighbour
search of the element: it keeps the first $30$ points found in element order within its search radius, so
near the surface, where the radius holds more points, the dropped points depend on the numbering. A rebuild of
the operator in Python on the same mesh shows the effect of this rule alone ($2\times10^{-3}$ to $0.2$
depending on the radius) and that a cap above the number of points, or a search that takes the nearest points,
removes it to round-off. All runs of this chapter keep the element's setting; [one $16^3$ run with a cap of
$120$ changes the main numbers by ..].
